\chapter{Proposed Methodology and Pipeline Architecture}
\label{ch:method}

\section{Overview}
The same training loop serves every task (Figure~\ref{fig:pipeline}). A set of Gaussians is rendered on the grid of
the unknown; a physics operator turns the rendered field into simulated measurements; the misfit to the real
measurements, plus an optional sparsity term, is back-propagated to all Gaussian parameters; and an adaptive density
control step periodically adds and removes Gaussians. For dynamic fields a deformation network maps the canonical
Gaussians to their state at each time step before rendering.

\begin{figure}[t]
\centering
\begin{tikzpicture}[
  box/.style={draw, rounded corners=3pt, align=center, minimum height=1.25cm, text width=2.55cm, font=\small, fill=blue!4},
  dark/.style={box, fill=blue!15},
  arr/.style={-{Stealth[length=2.5mm]}, thick},
  node distance=0.55cm]
  \node[box] (g) {Gaussians\\$\Theta=\{\mu_i,s_i,R_i,a_i\}$};
  \node[box, right=of g] (d) {Deformation\\$D_\phi(\mu_i,t)$\\{\footnotesize (4D only)}};
  \node[box, right=of d] (r) {Render on grid\\$f_\Theta(\vect{x})$\\{\footnotesize local window}};
  \node[box, right=of r] (a) {Physics\\$\mathcal{A}$\\{\footnotesize scattering / mask}};
  \node[dark, below=1.1cm of a] (l) {Loss\\$\mathcal{D}(\mathcal{A}f,\vect{y})+\lambda\,\mathcal{S}(f)$};
  \node[box, left=of l, text width=3.0cm] (y) {Measurements $\vect{y}$};
  \node[box, below=1.1cm of d, text width=3.4cm] (dc) {Densify / prune\\{\footnotesize at back-projected residual}};
  \draw[arr] (g) -- (d); \draw[arr] (d) -- (r); \draw[arr] (r) -- (a); \draw[arr] (a) -- (l); \draw[arr] (y) -- (l);
  \draw[arr, dashed] (l.south) -- ++(0,-0.5) -| node[pos=0.25, below, font=\footnotesize]{gradients (Adam)} (g.south);
  \draw[arr, dotted] (dc.north) -- ++(0,0.35) -| (g.south east);
\end{tikzpicture}
\caption{Pipeline. The camera rasteriser of 3DGS is replaced by rendering on a grid followed by the physics operator.}
\label{fig:pipeline}
\end{figure}

\section{Gaussian field representation}
Each Gaussian has a centre, per-axis scales, an orientation and an amplitude; the field at a grid point is the sum of
all Gaussians evaluated there (Section~\ref{sec:field}). Because a Gaussian is negligible beyond about three standard
deviations, each Gaussian is evaluated only in a window of $\pm h$ cells around its centre and accumulated into the
grid with a scatter-add, so the cost is $O(N(2h+1)^d)$ and independent of the grid size; scales are clamped to
$h/3$ so that the window always covers $\pm3\sigma$. For permittivity the amplitude is kept non-negative with a
softplus; for vorticity it is signed and the window wraps around the periodic domain. All parameters of one scene form
a single $N\times P$ tensor ($P=6$ in 2D, $P=11$ in 3D), which is convenient for sequence models.

\section{Physics operators}
\paragraph{Scattering.} We re-implement the InverseBench linear scattering operator as two complex matrix products
in single precision on the GPU and verify at start-up that it matches the benchmark implementation to within a relative
error of $3\times10^{-7}$. For 3D volumes, the operator is applied to every $z$-slice and the losses are summed.
\paragraph{Fluid observations.} For Navier--Stokes the measurement operator selects the observed pixels of each frame
(a random mask with 100\% or 10\% coverage); frames that are held out for evaluation are never observed.

\section{Training and adaptive density control}
All parameters are optimised with Adam~\citep{kingma2015adam} and a cosine learning-rate schedule. For scattering,
Gaussians are initialised uniformly at random, with no knowledge of the object. Every 250 iterations, during the
first 60\% of training, new Gaussians are inserted at the grid points with the largest magnitude of the
\emph{back-projected residual} $\mathrm{Re}\,\mathcal{A}^{\mathsf H}(\vect{y}-\mathcal{A}f)$, i.e.\ where the
measurements indicate missing structure, and Gaussians whose amplitude has fallen below $10^{-3}$ are pruned. This
replaces the image-space gradient criterion of 3DGS, which is unavailable without pixel supervision.

\section{Sparsity regulariser}
Early experiments showed that aggressive densification can place mass in regions to which the measurements are blind
(Section~\ref{sec:ghost}). A sparsity term on the total intensity of the rendered field penalises mass the data does
not require. Because the field is non-negative, the mean of the field equals its $\ell_1$ norm divided by the grid size.

\section{4D representation of fluid flow}
A canonical set of 2D Gaussians is initialised from the observed pixels of the first frame. A deformation network maps
each Gaussian's canonical centre and the time $t$ to offsets of all its parameters; positions move in pixel units.
Three deformation models are compared:
\begin{itemize}
  \item \textbf{MLP} (as in 4D Gaussian Splatting): Fourier features of $(\mu_i,t)$ fed to a three-layer MLP,
  evaluated independently for every Gaussian and time.
  \item \textbf{Mamba}: for each Gaussian the inputs at all $T$ time steps form a sequence processed by two Mamba
  blocks, so the offset at time $t$ depends on the Gaussian's earlier states (causal).
  \item \textbf{Bidirectional Mamba}: a second Mamba stack reads the sequence backwards and the two are summed, so
  every time step sees both past and future.
\end{itemize}
The output layer is initialised to zero so training starts from the canonical configuration. An optional temporal
smoothness term penalises differences between offsets at consecutive times, including unobserved times. A
\emph{static} model without deformation serves as an ablation.

\section{Baselines}
The baseline for scattering optimises every pixel (2D) or voxel (3D) directly with the same data term, starting from
zero and projected to $[0,1]$ after each step; optionally it uses the same sparsity term. Only the representation
differs between the two methods. For fluids, the static-Gaussian model is the baseline.

\section{Implementation}
The pipeline is a Python/PyTorch package (\texttt{thesis\_gs}) with modules for the Gaussian field, physics,
volume loading, scattering experiments, Navier--Stokes data, deformation networks and 4D training, a command-line
interface, experiment configurations and unit tests (Appendix~\ref{app:repro}). Mamba is implemented in pure PyTorch
with an explicit selective scan, which is sufficient for sequences of 21 steps and runs on any GPU.
